\documentclass[12pt]{article}
\usepackage[utf8]{inputenc}
\usepackage[T1]{fontenc}
\usepackage[margin=1in]{geometry}
\usepackage{setspace}
\usepackage{xcolor}
\usepackage[colorlinks=true, urlcolor=blue!60!black, linkcolor=blue!60!black]{hyperref}
\setlength{\parindent}{0em}
\setlength{\parskip}{0.45em}
\renewcommand{\baselinestretch}{1.0}
\begin{document}
\begin{center}
{\Large \textit{Ziglar v. Abbasi}, 582 U.S. 120 (2017)\par}
\vspace{0.3em}
{\normalsize Case Brief --- One-Page Edition\par}
\end{center}

\textbf{Memory Jogger.} No \textit{Bivens} suit against top officials over post-9/11 detention policy: expanding \textit{Bivens} is disfavored (\href{https://www.law.cornell.edu/supremecourt/text/15-1358}{Cornell LII}).

\textbf{Facts.} After 9/11, six men of Arab or South Asian descent were held for months in a Brooklyn maximum-security unit, allegedly confined 23 hours a day under constant light, strip-searched, and abused by guards with the wardens' knowledge. They sued Ashcroft, Mueller, Ziglar, and two wardens for damages.

\textbf{Procedural History.} The Second Circuit let most claims proceed against all defendants.

\textbf{Votes.} 4--2: Kennedy for the Court; Thomas concurred in part; Breyer, with Ginsburg, dissented; three Justices did not sit.

\textbf{Holding.} A claim differing meaningfully from \textit{Bivens}, \textit{Davis}, or \textit{Carlson} arises in a new context and fails if special factors counsel hesitation. The policy claims were barred, the abuse claim against Warden Hasty remanded, and the conspiracy claim met with qualified immunity.

\textbf{Reasoning.} \textit{Bivens}, \textit{Davis}, and \textit{Carlson} came from an ``ancien regime'' of implied remedies; Congress, not courts, should usually create them. New-context signs include the officers' rank, the right at issue, and the risk of intruding on other branches. High-level policy, national security, congressional silence despite oversight, and alternative remedies all counselled hesitation; ``even a modest extension is still an extension.'' Thomas would confine \textit{Bivens} to its facts; Breyer said the claims fit \textit{Bivens} and \textit{Carlson} and left detainees damages or nothing. \emph{Later:} \textit{Egbert} (2022) and \textit{Goldey} (2025).

\textbf{Examples.} \emph{Same-Side:} suing the DHS Secretary over emergency searches. \emph{Opposite-Side:} a warrantless home entry by a line agent. \emph{Fence-Sitter:} a pretrial detainee denied insulin.

\textbf{Critique.} Federal officers escape damages that state officers face under \S~1983, the alternative remedies are often illusory, and critics say \textit{Bivens} is being overruled in slow motion.

\textbf{Key Quotation.} ``The proper test for determining whether a case presents a new \textit{Bivens} context is as follows. If the case is different in a meaningful way from previous \textit{Bivens} cases decided by this Court, then the context is new.'' --- Justice Kennedy, 582 U.S. at 139.

\end{document}
